\documentclass[a4paper, serif, 10pt]{moderncv}

%% moderncv
% style and color
\moderncvstyle{banking}
\moderncvcolor{black}

%% page layout/margins
\usepackage[top=2.5cm, bottom=2.5cm, left=1.5cm, right=1.5cm]{geometry}

\nopagenumbers{}

%% Babel config for Indonesian (Bahasa Indonesia) language
% ref:
% - https://latex3.github.io/babel/guides/locale-indonesian.html
\usepackage[indonesian]{babel}

%% fontspec
\usepackage{fontspec}

\IfFontExistsTF{Linux Libertine}{
  \setmainfont[Ligatures={TeX, Common}, Numbers=OldStyle]{Linux Libertine}
}{}
\IfFontExistsTF{Linux Libertine O}{
  \setmainfont[Ligatures={TeX, Common}, Numbers=OldStyle]{Linux Libertine O}
}{}

%% CTeX
% CJK support, used to show CJK characters in the resume
%
% - fontset=none: disable builtin fontset but instead set the CJK font manually
% - heading=false: disable ctex heading
% - punct=kaiming: use kaiming punctuations styles for CJK
% - scheme=plain: use plain scheme, do not override \normalsize font size
% - space=auto: space settings for CJK characters
%
% ref:
% - http://ctan.mirrorcatalogs.com/language/chinese/ctex/ctex.pdf
\usepackage[UTF8, heading=false, punct=kaiming, scheme=plain, space=auto]{ctex}

\IfFontExistsTF{Noto Serif CJK SC}{
  \setCJKmainfont{Noto Serif CJK SC}
}{}
\IfFontExistsTF{Noto Sans CJK SC}{
  \setCJKsansfont{Noto Sans CJK SC}
}{}

%% line spacing
\usepackage{setspace}
\setstretch{1.125}

%% URL styling - use normal text instead of monospace
\urlstyle{same}

% Auto-underline all links
% Defer until \begin{document} so hyperref is already loaded
\AtBeginDocument{%
  \let\oldhref\href
  \renewcommand{\href}[2]{\oldhref{#1}{\underline{#2}}}%
}

%% Patch \httplink and \httpslink to handle full URLs
% The original moderncv macros always prepend http:// or https:// to URLs.
% This causes issues when the URL already has a protocol (e.g., https://example.com)
% resulting in malformed URLs like https://https://example.com.
% This patch checks if the URL already contains "://" and uses it directly if so.
% Note: We use \str_set:Nx (x-expansion) to expand macros like \@homepage first.
\ExplSyntaxOn
\str_new:N \g_yamlresume_url_str

\RenewDocumentCommand{\httpslink}{O{}m}{
  \str_set:Nx \g_yamlresume_url_str {#2}
  \str_if_in:NnTF \g_yamlresume_url_str {://}
    { \url{#2} }
    { \url{https://#2} }
}

\RenewDocumentCommand{\httplink}{O{}m}{
  \str_set:Nx \g_yamlresume_url_str {#2}
  \str_if_in:NnTF \g_yamlresume_url_str {://}
    { \url{#2} }
    { \url{http://#2} }
}
\ExplSyntaxOff

\name{Budi Santoso}{}
\title{Project Manager Berpengalaman di Bidang Teknologi Informasi}
\phone[mobile]{+62 812-3456-7890}
\email{budi.santoso@email.com}
\homepage{https://www.budisantoso.id}

\address{Jl. Jenderal Sudirman No. 52, Jakarta, DKI Jakarta, Indonesia, 12190}{}{}

\extrainfo{{\small \faLinkedin}\ \href{https://www.linkedin.com/in/budisantoso}{@budisantoso} {} {} {} • {} {} {} 
{\small \faGithub}\ \href{https://github.com/budisantoso}{@budisantoso}}

\begin{document}

\maketitle

\section{Informasi Dasar}

\cvline{}{\begin{itemize}
\item Project Manager dengan pengalaman lebih dari 8 tahun dalam memimpin proyek teknologi informasi dari perencanaan hingga penyelesaian.
\item Terampil dalam metodologi Agile, Scrum, dan Waterfall, serta memiliki rekam jejak sukses dalam mengelola tim lintas fungsi.
\item Mampu mengelola anggaran, jadwal, dan risiko proyek secara efektif untuk memastikan hasil yang sesuai dengan tujuan bisnis.
\item Berkomitmen pada pengembangan berkelanjutan dan komunikasi yang transparan dengan para pemangku kepentingan.
\end{itemize}}
\section{Pendidikan}

\cventry{Sep 2010–Jul 2014}
        {Sarjana, Ilmu Komputer, Nilai: 3.65}
        {Universitas Indonesia}
        {\url{https://www.ui.ac.id}}
        {}
        {\begin{itemize}
\item Memperoleh pemahaman mendalam tentang siklus hidup pengembangan perangkat lunak dan manajemen proyek.
\item Aktif dalam organisasi kemahasiswaan dan memimpin proyek akhir berbasis tim.
\end{itemize}
\textbf{Mata Kuliah}: Manajemen Proyek Perangkat Lunak, Rekayasa Perangkat Lunak, Sistem Basis Data, Jaringan Komputer, Algoritma dan Struktur Data}
\section{Pengalaman Kerja}

\cventry{Jan 2020–Sekarang}
        {Senior Project Manager}
        {PT Teknologi Maju}
        {\url{https://www.teknologimaju.co.id}}
        {}
        {\begin{itemize}
\item Memimpin proyek pengembangan platform digital berskala enterprise dengan anggaran lebih dari 10 miliar rupiah.
\item Mengelola tim berisi 20+ anggota lintas fungsi, termasuk developer, desainer, dan QA.
\item Menerapkan kerangka kerja Agile dan Scrum untuk meningkatkan kecepatan rilis produk hingga 30\%.
\item Berhasil menyelesaikan 15+ proyek tepat waktu dengan tingkat kepuasan pelanggan 95\%.
\end{itemize}
\textbf{Kata Kunci}: Agile, Scrum, Manajemen Proyek, Kepemimpinan, Enterprise}

\cventry{Mar 2017–Des 2019}
        {Project Manager}
        {PT Digital Solusi}
        {\url{https://www.digitalsolusi.co.id}}
        {}
        {\begin{itemize}
\item Mengelola proyek pengembangan aplikasi mobile dan web untuk klien di sektor perbankan dan ritel.
\item Menyusun rencana proyek, jadwal, dan anggaran, serta memantau kemajuan melalui tools seperti JIRA dan Trello.
\item Berkoordinasi dengan tim internal dan eksternal untuk memastikan kebutuhan klien terpenuhi.
\item Meningkatkan efisiensi proses kerja tim melalui adopsi Kanban dan otomatisasi pelaporan.
\end{itemize}
\textbf{Kata Kunci}: Manajemen Proyek, JIRA, Trello, Mobile, Web}

\cventry{Agu 2014–Feb 2017}
        {Junior Project Manager}
        {PT Inovasi Kreatif}
        {\url{https://www.inovasikreatif.co.id}}
        {}
        {\begin{itemize}
\item Membantu manajer proyek senior dalam perencanaan dan eksekusi proyek teknologi.
\item Mengelola dokumentasi proyek, laporan status, dan komunikasi dengan pemangku kepentingan.
\item Berkontribusi pada pengembangan metodologi internal untuk manajemen proyek yang lebih terstruktur.
\end{itemize}
\textbf{Kata Kunci}: Dokumentasi, Laporan, Komunikasi, Waterfall}
\section{Bahasa}

\cvline{Indonesia}{Bahasa Ibu / Bilingual \hfill \textbf{Kata Kunci}: Bahasa Ibu}
\cvline{Inggris}{Tingkat Profesional Penuh \hfill \textbf{Kata Kunci}: TOEFL 110, IELTS 7.5}
\section{Keahlian}

\cvline{Manajemen Proyek}{Ahli \hfill \textbf{Kata Kunci}: Agile, Scrum, Kanban, Waterfall, Manajemen Risiko, Anggaran}
\cvline{Tools}{Mahir \hfill \textbf{Kata Kunci}: JIRA, Trello, MS Project, Asana, Confluence, Slack}
\cvline{Kepemimpinan}{Mahir \hfill \textbf{Kata Kunci}: Motivasi Tim, Delegasi, Negosiasi, Komunikasi, Problem Solving}
\section{Penghargaan}

\cventry{Des 2022}
        {PT Teknologi Maju}
        {Project Manager Terbaik}
        {}
        {}
        {Diberikan sebagai pengakuan atas kepemimpinan luar biasa dalam menyelesaikan proyek strategis perusahaan dengan hasil melebihi target.}

\cventry{Nov 2019}
        {PT Digital Solusi}
        {Inovasi Proses Bisnis}
        {}
        {}
        {Penghargaan atas kontribusi dalam merancang ulang proses manajemen proyek yang meningkatkan efisiensi tim sebesar 25\%.}
\section{Sertifikat}

\cventry{Jun 2021}
        {Project Management Institute}
        {Project Management Professional (PMP)}
        {\url{https://www.pmi.org/certifications/project-management-pmp}}
        {}
        {}

\cventry{Mar 2019}
        {Scrum Alliance}
        {Certified ScrumMaster (CSM)}
        {\url{https://www.scrumalliance.org/certifications}}
        {}
        {}
\section{Publikasi}

\cventry{Agu 2020}
        {Meningkatkan Keberhasilan Proyek TI melalui Kombinasi Agile dan Waterfall}
        {Jurnal Manajemen Proyek Indonesia}
        {\url{https://www.jmpi.id/artikel/agile-waterfall}}
        {}
        {\begin{itemize}
\item Artikel ini membahas pendekatan hybrid dalam manajemen proyek teknologi informasi.
\item Menyajikan studi kasus dari beberapa proyek enterprise di Indonesia.
\item Memberikan rekomendasi praktis untuk manajer proyek dalam memilih metodologi yang tepat.
\end{itemize}}
\section{Proyek}

\cventry{Jan 2021–Des 2021}
        {Pengembangan platform e-commerce untuk perusahaan ritel besar di Indonesia}
        {Platform E-Commerce Terintegrasi}
        {\url{https://www.teknologimaju.co.id/proyek/ecommerce}}
        {}
        {\begin{itemize}
\item Memimpin tim berisi 15 orang dalam membangun platform e-commerce dari awal.
\item Mengimplementasikan arsitektur microservices dan integrasi dengan sistem pembayaran.
\item Mencapai waktu muat halaman di bawah 2 detik dan meningkatkan konversi penjualan sebesar 40\%.
\end{itemize}
\textbf{Kata Kunci}: E-Commerce, Microservices, Agile}

\cventry{Mar 2018–Nov 2019}
        {Sistem manajemen rumah sakit untuk digitalisasi rekam medis dan operasional}
        {Sistem Informasi Rumah Sakit}
        {\url{https://www.digitalsolusi.co.id/proyek/simrs}}
        {}
        {\begin{itemize}
\item Mengelola proyek pengembangan SIMRS untuk tiga rumah sakit di Jakarta.
\item Berkoordinasi dengan tenaga medis dan staf administrasi untuk mengumpulkan kebutuhan.
\item Sistem berhasil mengurangi waktu pemrosesan rekam medis hingga 50\%.
\end{itemize}
\textbf{Kata Kunci}: Kesehatan, SIMRS, Waterfall}
\section{Minat}

\cvline{Olahraga}{Lari, Renang, Bersepeda}
\cvline{Teknologi}{AI, Blockchain, IoT}
\section{Kegiatan Sukarela}

\cventry{Jan 2018–Des 2022}
        {Mentor Manajemen Proyek}
        {Yayasan Pendidikan Teknologi}
        {\url{https://www.ypt.or.id}}
        {}
        {\begin{itemize}
\item Menjadi mentor bagi mahasiswa dan profesional muda dalam bidang manajemen proyek teknologi.
\item Mengadakan workshop tentang Agile, Scrum, dan tools manajemen proyek.
\item Membantu peserta mengembangkan keterampilan kepemimpinan dan komunikasi.
\end{itemize}}
    
\end{document}